\section{Study design}
\label{sec:design}

\subsection{Data and evaluation protocol}

We use the human breast carcinoma subset of MIDOG++ \citep{aubreville2023comprehensive}: 150 cases, 50 from each of
three scanners (Hamamatsu XR, Hamamatsu S360, Aperio CS2), with 1{,}721 annotated mitotic figures and 2{,}714 hard
negatives (imposters), i.e.\ 4{,}435 cells. Each model classifies a patch centred on an annotation (mitotic figure vs.\
imposter); candidate detection is outside our scope. We hold out one scanner at a time and train on the official training images of the other two, except eight images per training scanner that
form a fixed validation set. Checkpoint and decision threshold (\emph{cutoff}) are always selected on this source
validation set by F1, so the held-out scanner never informs a training decision.

All 50 images of the held-out scanner are used for evaluation through cross-fitting. They are sorted by
identifier and split into alternating halves $A$ and $B$. Every method that uses unlabelled new-scanner data is fitted
on one half and applied to the other, then the halves are swapped, so each cell is scored by a model that never saw its
image in any form. The outcome is F1 on all cells of the held-out scanner, computed per run, averaged over the student
seeds and then over the three held-out scanners. New-scanner labels are used only for scoring and for the
labelled-cutoff upper bound.

\subsection{Teachers and students}

\emph{Teachers.} Four pathology foundation models, H-optimus-0 \citep{hoptimus0}, Virchow2 \citep{zimmermann2024virchow2},
UNI \citep{chen2024uni} and UNI2-h \citep{uni2h}, are adapted with LoRA \citep{hu2022lora} (rank 16, $\alpha=16$,
dropout 0.1 on every attention and MLP projection; frozen encoder; linear head on the class token concatenated with the
mean patch token), following a recent mitotic-figure benchmark \citep{ammeling2026benchmarking}. Each teacher is
trained for 5 epochs on 224-pixel centre crops, once per held-out scanner and LoRA seed (three seeds).

\emph{Students.} To test whether the decomposition depends on BatchNorm, we use two ImageNet-initialised students:
ResNet-18 \citep{he2016deep} (11\,M parameters, BatchNorm) and ConvNeXt-T \citep{liu2022convnet} (28\,M parameters,
LayerNorm, no statistics to re-estimate), trained on 128-pixel crops for 60 epochs with 10 seeds per configuration. The
\emph{no-teacher student} is trained with expert labels only. The \emph{distilled (KD) student}
\citep{hinton2015distilling}, in contrast, never sees an expert label: it matches the teacher's temperature-softened
probability ($T=2$) and, through a linear projection, the teacher's feature vector (cosine loss, weight 1) on 16 cached
views of each source cell (8 dihedral transforms, each with and without H\&E stain jitter). Students are distilled from
the LoRA seed-0 teacher unless stated otherwise.

\subsection{Three layers and their repairs}
\label{sec:layers}

We describe what the student loses on a new scanner in three layers and repair them in order of cost
(\figureref{fig:schematic}; \tableref{tab:ladder} in Appendix~\ref{app:impl}).

\emph{Operating point.} Even when a student keeps its ranking (AUROC) on a new scanner, a cutoff chosen on source data
can be wrong there \citep{roschewitz2023automatic,thiringer2026scanner}. Our repair, \emph{teacher-referenced
recalibration} (TGR), keeps the student's weights and fits one logit offset so that the student agrees with the teacher
on the unlabelled cells $U$ of the fit half:
\begin{equation}
\hat b = \arg\min_b \; -\frac{1}{|U|}\sum_{x\in U}\Big[q(x)\log\sigma\big(z(x)+b\big) + \big(1-q(x)\big)\log\big(1-\sigma(z(x)+b)\big)\Big],
\label{eq:tgr}
\end{equation}
where $z(x)$ is the student logit and $q(x)$ the teacher probability; a cell of the other half is called positive if
$z(x)+\hat b \ge \operatorname{logit}(\tau_T)$, with $\tau_T$ the teacher's own source cutoff. Because the ranking is
kept, this shift carries the teacher's view of the target prevalence over to the student. TGR needs no retraining, only the
teacher's forward pass on a few unlabelled images (about 150 cells for five).

\emph{Normalisation statistics.} For ResNet-18, AdaBN \citep{li2018adaptive,schneider2020improving} re-estimates the
BatchNorm running statistics on the fit half without changing any learned weight. TENT \citep{wang2021tent} additionally
tunes the BatchNorm affine parameters by entropy minimisation (learning rate $10^{-3}$, one epoch; Appendix~\ref{app:tent}).

\emph{Representation.} Teacher-guided adaptation (TGA) retrains the KD student with the unlabelled cells of the fit
half added to the distillation set, using the teacher's answers and features on the same views. As a control for the
contribution of the teacher, \emph{no-teacher adaptation} trains on the same
unlabelled cells with the soft answers of an ensemble of the ten no-teacher ResNet-18 students after AdaBN on the same
half, and with expert labels on the source cells; no foundation model is involved, and the same pseudo-labeller serves
both students.

\subsection{Attribution}

With $F_\text{ctrl}$ the new-scanner F1 of the no-teacher student and $F_\text{teacher}$ that of the teacher at its
source cutoff, we express every arm by the \emph{share of the teacher's advantage recovered},
$(F_\text{arm}-F_\text{ctrl})/(F_\text{teacher}-F_\text{ctrl})$ (0\% = no-teacher student, 100\% = teacher), the
``performance gap recovered'' of \citet{burns2023weak}. Because AdaBN also moves the operating point, the credit of the
two post-hoc layers depends on their order. We therefore report Shapley values over \{threshold, statistics\}
\citep{zhang2023why}, i.e.\ each layer's gain averaged over the two orders. Retraining is credited with the gain of TGA beyond the two post-hoc layers, and the rest is the remaining gap.

\subsection{Statistical analysis}
\label{sec:stats}

Uncertainty comes from a paired bootstrap \citep{efron1993introduction} with $B=20{,}000$ replicates that resamples the
50 images of each held-out scanner (multinomial weights on images, shared by all arms) and the 10 student seeds, since
both are sources of variation \citep{varoquaux2022machine}; intervals are 95\% percentile intervals of the
fold-averaged statistic. Two-sided bootstrap $p$-values are adjusted with Holm's procedure \citep{holm1979simple} within
each family of contrasts.

Each family of confirmatory contrasts, together with its expected direction, was written down and its file hash
recorded before the corresponding models were trained or the corresponding numbers computed (Appendix~\ref{app:prereg});
results that contradict an expectation are reported as such, and analyses added afterwards are labelled descriptive. To
ask whether the student is as good as its teacher, we test non-inferiority with a margin of $\delta=0.02$ F1, the range
of the H-optimus-0 teacher's own new-scanner F1 over three LoRA seeds (0.801 to 0.821): the student is called
non-inferior if the lower limit of the 95\% interval of (student $-$ teacher) exceeds $-0.02$ (one-sided $\alpha=0.025$,
Holm within the family). A failed test is reported as ``non-inferiority not established'', not as equivalence or
difference.
